\itemhead{习题1.2 · 第9题（教材p.24）}{XJ-B1-C01-HW-P024-Y1370}
\label{XJ-B1-C01-HW-P024-Y1370}
{\footnotesize 湘教版教材p.24 / PDF第31页 / 习题1.2 / 第9题}
\field{原题}{对下列含有量词的命题作否定，并判断其真假：
\subq{1}$p$：所有能被 2 整除的数都是偶数；
\subq{2}$p$：$\exists x_0\in\mathbb N,\ 2^{x_0}<0$；
\subq{3}$p$：$\exists x_0\in\complement_{\mathbb R}\mathbb Q,\ x_0^2\in\mathbb Q$；
\subq{4}$p$：$\exists x\in\mathbb Z_+,\ \dfrac{3x+2}{x+1}\in\mathbb N$。}
\field{参考答案}{（1）存在能被2整除却不是偶数的整数，假；（2）$\forall x_0\in\mathbb N,2^{x_0}\ge0$，真；（3）$\forall x_0\in\mathbb R\setminus\mathbb Q,x_0^2\notin\mathbb Q$，假（取 $\sqrt2$）；（4）$\forall x\in\mathbb Z_+,(3x+2)/(x+1)\notin\mathbb N$，真，因为该式为 $3-1/(x+1)$。}
{\small \field{关联知识}{L04 含一个量词命题的否定（教材帮p.39、40 / PDF48、49）}
\field{方法与易错点}{分三步书写：保留论域、交换量词、否定断言，再判断新命题真假。 易错：只改量词不改结论；把不等号否定成反向但遗漏等号。}
\field{答案依据}{答案册相关题组：PDF第7页（册内页6）；按当前题库原题逐项重算，不按旧版题号机械对套。}}
